\chapter{Single-sorted underlying transformations}
\label{ch:ss_underlying}

Every single-sorted $n$-category contains a lower-dimensional $m$-category obtained by restricting
to the $m$-cells: one can show that the source, target, and composition operations preserve the set
of $m$-cells, so we may consider their birestrictions to it. This construction extends naturally to
the $\omega$-categorical setting. Moreover, any single-sorted $n$-functor maps $m$-cells to
$m$-cells, so it restricts to a functor between the resulting lower-dimensional categories. We call
this process the \emph{underlying} transformation, and we show that it is functorial and that
underlying transformations at different dimensions are compatible with one another.

We conclude the chapter by showing that the categories $n\Cat$ together with the underlying functors
form a projective system, and that $\omega\Cat$ is its projective limit.

\section{Underlying single-sorted categories}

We begin by showing that, for any dimension $k < m$, the source, target, and composition operations
at dimension $k$ of a single-sorted category can be birestricted to the set of $m$-cells.

\begin{lemma}
  \label{lem:underlying_source_is_cell}
  \uses{def:cell}
  \lean{HigherCategoryTheory.SingleSorted.underlying_source_is_cell}
  \leanok
  Let $\cat{C}$ be a single-sorted $\Index$-category and $m \in \Index$. Then, for every $f \in
  \cat{C}$ and every $k \in \Index$ such that $k < m$, we have that $\Sc^k (f)$ is a $m$-cell in
  $\cat{C}$.
\end{lemma}

\begin{proof}
  \leanok
  Applying one of the axioms of single-sorted categories, we have that
  \[
    \Sc^m (\Sc^k f) = \Sc^k (f).
  \]
\end{proof}

\begin{corollary}
  \label{cor:underlying_source_birestriction}
  \uses{lem:underlying_source_is_cell}
  Let $\cat{C}$ be a single-sorted $\Index$-category and $m \in \Index$. Then, for every $k \in
  \Index$ such that $k < m$, the source operation $\Sc^k$ birestricts to the set of $m$-cells $C_m$.
\end{corollary}

\begin{proof}
  Follows directly from Lemma~\ref{lem:underlying_source_is_cell}.
\end{proof}

\begin{lemma}
  \label{lem:underlying_target_is_cell}
  \uses{thm:cell_sc_iff_cell_tg}
  \lean{HigherCategoryTheory.SingleSorted.underlying_target_is_cell}
  \leanok
  Let $\cat{C}$ be a single-sorted $\Index$-category and $m \in \Index$. Then, for every $f \in
  \cat{C}$ and every $k \in \Index$ such that $k < m$, we have that $\Tg^k (f)$ is a $m$-cell in
  $\cat{C}$.
\end{lemma}

\begin{proof}
  \leanok
  Applying one of the axioms of single-sorted categories, we have that
  \[
    \Tg^m (\Tg^k f) = \Tg^k (f).
  \]
  That is, $\Tg^k (f)$ is a target-based $m$-cell in $\cat{C}$. However, we know that every
  target-based $m$-cell is also a $m$-cell.
\end{proof}

\begin{corollary}
  \label{cor:underlying_target_birestriction}
  \uses{lem:underlying_target_is_cell}
  Let $\cat{C}$ be a single-sorted $\Index$-category and $m \in \Index$. Then, for every $k \in
  \Index$ such that $k < m$, the target operation $\Tg^k$ birestricts to the set of $m$-cells $C_m$.
\end{corollary}

\begin{proof}
  Follows directly from Lemma~\ref{lem:underlying_target_is_cell}.
\end{proof}

\begin{lemma}
  \label{lem:underlying_comp_is_cell}
  \uses{def:cell}
  \lean{HigherCategoryTheory.SingleSorted.underlying_comp_is_cell}
  \leanok
  Let $\cat{C}$ be a single-sorted $\Index$-category and $m \in \Index$. Then, for every $f, g \in
  C_m$ and every $k \in \Index$ such that $k < m$, if $f$ and $g$ are composable at dimension $k$,
  then $g \Pcomp^k f$ is an $m$-cell in $\cat{C}$.
\end{lemma}

\begin{proof}
  \leanok
  Applying one of the axioms of single-sorted categories and the fact that $f$ and $g$ are
  $m$-cells, we have that
  \[
    \Sc^m (g \Pcomp^k f) = \Sc^m (g) \Pcomp^k \Sc^m (f) = g \Pcomp^k f.
  \]
\end{proof}

\begin{corollary}
  \label{cor:underlying_comp_birestriction}
  \uses{lem:underlying_comp_is_cell}
  Let $\cat{C}$ be a single-sorted $\Index$-category and $m \in \Index$. Then, for every $k \in
  \Index$ such that $k < m$, the composition operation $\Pcomp^k$ birestricts to the set of
  $m$-cells $C_m$.
\end{corollary}

\begin{proof}
  Follows directly from Lemma~\ref{lem:underlying_comp_is_cell}.
\end{proof}

\begin{construction}[Underlying single-sorted finite category]
  \label{def:NCategory_underlying}
  \uses{
    cor:underlying_source_birestriction,
    cor:underlying_target_birestriction,
    cor:underlying_comp_birestriction
  }
  \lean{HigherCategoryTheory.SingleSorted.NCategory.underlying}
  \notready
  Let $\cat{C}$ be a single-sorted $n$-category, and let $m \in \NN$ with $m < n$. The
  \emph{underlying $m$-category} of $\cat{C}$ is the single-sorted $m$-category $\cat{C}^{<m}$
  defined on the set of $m$-cells $C_m$, with the source, target, and composition operations given
  by the birestrictions of those of $\cat{C}$ to $C_m$.
\end{construction}

\begin{proof}
  The birestrictions are well-defined by Corollaries~\ref{cor:underlying_source_birestriction},
  \ref{cor:underlying_target_birestriction}, and~\ref{cor:underlying_comp_birestriction}. Each axiom
  of single-sorted categories for $\cat{C}^{<m}$ follows directly from the corresponding axiom for
  $\cat{C}$.
\end{proof}

\begin{construction}[Underlying single-sorted omega category]
  \label{def:OmegaCategory_underlying}
  \uses{
    def:SingleSortedOmegaCategory,
    cor:underlying_source_birestriction,
    cor:underlying_target_birestriction,
    cor:underlying_comp_birestriction
  }
  \lean{HigherCategoryTheory.SingleSorted.OmegaCategory.underlying}
  \notready
  Let $\cat{C}$ be a single-sorted $\omega$-category, and let $m \in \NN$. The \emph{underlying
  $m$-category} of $\cat{C}$ is the single-sorted $m$-category $\cat{C}^{<m}$ defined by the same
  construction as in Definition~\ref{def:NCategory_underlying}.
\end{construction}

\begin{proof}
  The single-sorted category axioms follow by the same argument as in
  Definition~\ref{def:NCategory_underlying}.
\end{proof}

\section{Underlying single-sorted functors}

\begin{lemma}
  \label{lem:underlying_functor_is_cell}
  \uses{def:cell, def:SingleSortedFunctor}
  \lean{HigherCategoryTheory.SingleSorted.underlying_functor_is_cell}
  \leanok
  Let $\cat{C}$ and $\cat{D}$ be single-sorted $\Index$-categories, $F \colon \cat{C} \to \cat{D}$
  be a single-sorted functor, and $m \in \Index$. Then, for every $f \in C_m$ we have that $F(f) \in
  D_m$.
\end{lemma}

\begin{proof}
  \leanok
  Applying one of the axioms of single-sorted functors and the fact that $f$ is an $m$-cell, we have
  that
  \[
    \Sc^m (F(f)) = F(\Sc^m (f)) = F(f).
  \]
\end{proof}

\begin{corollary}
  \label{cor:underlying_functor_birestriction}
  \uses{lem:underlying_functor_is_cell}
  Let $\cat{C}$ and $\cat{D}$ be single-sorted $\Index$-categories, $F \colon \cat{C} \to \cat{D}$
  be a single-sorted functor, and $m \in \Index$. Then $F$ birestricts to a function from $C_m$ to
  $D_m$.
\end{corollary}

\begin{proof}
  Follows directly from Lemma~\ref{lem:underlying_functor_is_cell}.
\end{proof}

\begin{construction}[Underlying single-sorted finite functor]
  \label{def:NFunctor_underlying}
  \uses{def:NCategory_underlying, cor:underlying_functor_birestriction}
  \lean{HigherCategoryTheory.SingleSorted.NFunctor.underlying}
  \notready
  Let $\cat{C}$ and $\cat{D}$ be single-sorted $n$-categories, and let $F \colon \cat{C} \to
  \cat{D}$ be a single-sorted $n$-functor. For any $m \in \NN$ with $m < n$, the birestriction of
  $F$ to the sets of $m$-cells defines a single-sorted $m$-functor $F^{<m} \colon \cat{C}^{<m} \to
  \cat{D}^{<m}$.
\end{construction}

\begin{proof}
  The birestriction is well-defined by Corollary~\ref{cor:underlying_functor_birestriction}. Each
  preservation axiom for $F^{<m}$ follows directly from the corresponding axiom for $F$.
\end{proof}

\begin{construction}[Underlying single-sorted omega functor]
  \label{def:OmegaFunctor_underlying}
  \uses{def:OmegaCategory_underlying, cor:underlying_functor_birestriction}
  \lean{HigherCategoryTheory.SingleSorted.OmegaFunctor.underlying}
  \notready
  Let $\cat{C}$ and $\cat{D}$ be single-sorted $\omega$-categories, and let $F \colon \cat{C} \to
  \cat{D}$ be a single-sorted $\omega$-functor. For any $m \in \NN$, the birestriction of $F$ to the
  sets of $m$-cells defines a single-sorted $m$-functor $F^{<m} \colon \cat{C}^{<m} \to
  \cat{D}^{<m}$.
\end{construction}

\begin{proof}
  The argument is identical to that of Definition~\ref{def:NFunctor_underlying}.
\end{proof}

\section{Single-sorted underlying functors}

The underlying constructions of the previous sections are in fact functorial, as the following
definition makes precise.

\begin{construction}[Underlying functor]
  \label{def:UnderlyingFunctor}
  \uses{
    def:NCategory_underlying,
    def:OmegaCategory_underlying,
    def:NFunctor_underlying,
    def:OmegaFunctor_underlying
  }
  \lean{HigherCategoryTheory.SingleSorted.UnderlyingFunctor}
  \notready
  Let $n, m \in \NInf$ with $m \leq n$. The \emph{underlying functor} $\Underlying^{(n, m)} \colon
  n\Cat \to m\Cat$ is the functor that maps every single-sorted $n$-category $\cat{C}$ to its
  underlying $m$-category $\cat{C}^{<m}$ and every single-sorted $n$-functor $F \colon \cat{C} \to
  \cat{D}$ to its underlying $m$-functor $F^{<m} \colon \cat{C}^{<m} \to \cat{D}^{<m}$. If $n = m$,
  then $\Underlying^{(n, n)}$ stands for the identity functor on $n\Cat$.
\end{construction}

\begin{proof}
  Functoriality follows from the fact that the underlying constructions preserve the birestriction
  structure: the identity functor on $\cat{C}$ yields the identity functor on $\cat{C}^{<m}$, and
  the composition of two $n$-functors yields the composition of their underlying $m$-functors.
\end{proof}

Moreover, underlying transformations at different dimensions are compatible with one another.

\begin{theorem}[Compatibility of underlying functors]
  \label{thm:UnderlyingFunctor_compatibility}
  \uses{def:UnderlyingFunctor}
  \notready
  The assignment $(n, m) \mapsto \Underlying^{(n, m)}$ defines a functor from $(\NInf, \leq)^{op}$
  to $\cat{Cat}$. That is, for all $k \leq m \leq n$ in $\NInf$,
  \[
    \Underlying^{(m, k)} \circ \Underlying^{(n, m)} = \Underlying^{(n, k)}.
  \]
\end{theorem}

\begin{proof}
  Both sides act as the same birestriction on the underlying map. The categorical structures
  coincide because restricting to $m$-cells and then to $k$-cells is the same as restricting to
  $k$-cells directly, since every $k$-cell is also an $m$-cell.
\end{proof}

\section{The projective limit of single-sorted categories}

\begin{theorem}[Projective limit of single-sorted categories]
  \label{thm:Underlying_projective_limit}
  \uses{thm:UnderlyingFunctor_compatibility}
  \notready
  The category $\omega\Cat$ is the projective limit of the system $(n\Cat, \Underlying^{(n, m)})_{m
  \leq n}$.
\end{theorem}
